\documentclass[conference]{IEEEtran}
\usepackage{booktabs}
\usepackage{amsmath,amssymb}
\usepackage{array}
\usepackage{tabularx}
\usepackage{url}
\usepackage{microtype}
\usepackage{balance}
\emergencystretch=2em

\begin{document}
\title{GovRed-Health: Benchmarking Authorization Drift and Bitemporal Governance Between Disclosure and Persistent Commit}
\author{\IEEEauthorblockN{Nguyen Ngoc Thien, Trinh Minh Quang}
\IEEEauthorblockA{Hai Ba Trung High School, Hue, Vietnam\\
Corresponding author: kakangocthien109@gmail.com}}
\maketitle

\begin{abstract}
Stateful clinical AI systems operating over longitudinal healthcare big data repeatedly query historical patient context and subsequently emit persistent state modifications. However, in distributed clinical workflows, patient state, bitemporal valid/knowledge horizons, dynamic consent, governing policy, role permissions, purpose constraints, cache state, or snapshot validity can drift across the disclosure-to-commit interval ($\Delta t = t_{\text{commit}} - t_{\text{disc}}$). Rooted in Snodgrass's foundational bitemporal data systems framework, GovRed-Health establishes an empirical benchmark for authorization drift in stateful clinical big data environments, combining a controlled logical schedule suite ($N=270$ schedules per arm) with an extensive retrospective clinical trace benchmark across 10,000+ real-world hospital event sequences from the MIMIC-IV database. In controlled execution, GLHS-Strict reduced the non-safe composite from 57.1\% (Unbound) to 14.3\% ($p=1.62\times 10^{-27}$), with zero serial invalid commits and zero sentinel PHI leakage (0/270). In retrospective MIMIC-IV workflow replay, unminimized agents suffered a 12.4\% empirical TOCTOU drift rate due to intervening laboratory and medication order modifications; GLHS-Strict eliminated this drift completely (0.00\% invalid writes) while reducing prompt token volume by 87.4\% and preserving complete bitemporal audit lineage.
\end{abstract}

\begin{IEEEkeywords}
healthcare big data, bitemporal data systems, health AI, authorization drift, TOCTOU, longitudinal governance, zero-PHI leakage, persistent agents, benchmark
\end{IEEEkeywords}

\section{Introduction}
Longitudinal healthcare big data platforms aggregate multi-source, append-only clinical event streams across electronic health records (EHRs), remote telemetry, and diagnostic repositories. In stateful clinical AI architectures, language model agents repeatedly read from these longitudinal patient stores to formulate therapeutic recommendations, draft clinical summaries, or trigger persistent state changes. This decoupled interaction model introduces a fundamental cross-operation integrity vulnerability: a governed disclosure can be fully valid when presented to a model at inference time $t_{\text{disc}}$, yet become an unauthorized or invalid basis for a durable write at commit time $t_{\text{commit}}$ after patient state, clinical evidence, dynamic consent, institutional policy, clinician role, or snapshot validity has changed.

This authorization-drift phenomenon represents a Time-of-Check to Time-of-Use (TOCTOU) problem that is distinct from single-turn prompt injection or isolated retrieval errors. Existing evaluation frameworks focus primarily on adversarial jailbreaks (AgentDojo~\cite{agentdojo}, ASB~\cite{asb}), medical prompt injection (MPIB~\cite{mpib}), or sensitive-read confinement (SecureClaw~\cite{secureclaw}). In parallel, agent transaction research has advanced commit-time authority checking (CommitGuard~\cite{commitguard}), policy-state serializability (Provenact~\cite{provenact}), multi-principal memory access control (GateMem~\cite{gatemem}), and bitemporal memory algebras (TOKI~\cite{toki}). As summarized in Table~\ref{tab:benchmark_comparison}, prior benchmarks evaluate isolated prompt or access-control failures but do not capture the bitemporal authorization drift that occurs across the disclosure-to-commit interval.

\begin{table}[t]
\caption{Cross-Benchmark Capability Comparison: Threat Dimension and System Mechanism Coverage in Stateful Health AI Security.}
\label{tab:benchmark_comparison}
\centering\scriptsize
\setlength{\tabcolsep}{3pt}
\begin{tabularx}{\columnwidth}{l c c c c c c}
\toprule
\textbf{Benchmark} & \textbf{Prompt Attack} & \textbf{Tool Access} & \textbf{TOCTOU Drift} & \textbf{Consent Revoke} & \textbf{Bitemporal} & \textbf{Merkle Trace}\\
\midrule
AgentDojo~\cite{agentdojo} & \checkmark & \checkmark & $\times$ & $\times$ & $\times$ & $\times$\\
MPIB~\cite{mpib} & \checkmark & $\times$ & $\times$ & $\times$ & $\times$ & $\times$\\
CommitGuard~\cite{commitguard} & $\times$ & \checkmark & \checkmark & $\times$ & $\times$ & $\times$\\
Provenact~\cite{provenact} & $\times$ & \checkmark & \checkmark & $\times$ & $\times$ & $\times$\\
GateMem~\cite{gatemem} & $\times$ & \checkmark & $\times$ & $\times$ & $\times$ & $\times$\\
TOKI~\cite{toki} & $\times$ & \checkmark & $\times$ & $\times$ & \checkmark & $\times$\\
\textbf{GovRed-Health (Ours)} & \checkmark & \checkmark & \textbf{\checkmark} & \textbf{\checkmark} & \textbf{\checkmark} & \textbf{\checkmark}\\
\bottomrule
\end{tabularx}
\end{table}

GovRed-Health addresses the specific systems-governance gap in healthcare big data: \emph{what happens after a governed, task-bounded health disclosure has been consumed by an AI agent and the underlying longitudinal state, bitemporal timeline, or governance authority mutates prior to persistent durability?} The benchmark incorporates Richard T. Snodgrass's foundational bitemporal data systems theory~\cite{snodgrass1995,kumar1998}, distinguishing valid time (when a clinical assertion was factually true) from transaction/knowledge time (when it was physically committed into the database). By embedding bitemporal coordinates and cryptographic Merkle lease seals into the disclosure object, GovRed-Health benchmarks scalable longitudinal governance with empirical zero-PHI leakage across multi-tenant persistence layers.

Its contributions are: (1) a frozen schedule abstraction incorporating model-visible bitemporal disclosures into the statistical unit; (2) matched state and governance mutations spanning subject, consent, policy, role/purpose, bitemporal state, cache, and audit dimensions; (3) end-to-end execution across PostgreSQL, Redis, and cryptographic audit persistence boundaries; and (4) an evaluation protocol distinguishing safe defensive rejection, transport failure, and indeterminate concurrent linearization under zero-PHI leakage guarantees.

\section{Bitemporal Governance and Threat Model}
\subsection{What GovRed-Health does not claim}
GovRed-Health does not introduce prompt injection, access-control theory, bitemporal database algebra, dual-boundary enforcement, optimistic concurrency control, or commit-time authorization. The four evaluated system arms are internal mechanism ablations designed to isolate individual architectural components, not claimed reference implementations of prior published systems.

\subsection{Snodgrass Bitemporal Foundations}
Clinical big data repositories cannot rely on a single physical timestamp. Following Snodgrass's temporal database foundations~\cite{snodgrass1995}, clinical assertions $\mathcal{E}$ are characterized by two orthogonal time dimensions:
\begin{enumerate}
\item \textbf{Valid Time} ($\tau_v = [v_{\text{start}}, v_{\text{end}})$): The continuous time interval during which a physiological state, clinical condition, or medication regimen was active in the patient's biological reality.
\item \textbf{Transaction / Knowledge Time} ($\tau_k = [t_{\text{start}}, t_{\text{end}})$): The discrete time interval during which the assertion was physically recorded, known to the clinical data store, and unmodified.
\end{enumerate}
Retrospective clinical chart amendments, backdated diagnostic entries, or late-arriving laboratory results alter $\tau_k$ or $\tau_v$. When an AI agent performs multi-step longitudinal reasoning across $\Delta t = t_{\text{commit}} - t_{\text{disc}}$, retrospective changes create bitemporal drift, rendering the agent's inferential premise invalid even if current point-in-time state appears superficially fresh.

\subsection{Authorization Drift Object}
We define an issued governed disclosure snapshot as a bitemporal tuple:
\begin{equation}
H = (u, a, p, \tau_v, \tau_k, v_s, v_g, d, e, E),
\end{equation}
where $u$ is the patient subject, $a$ the authenticated clinical actor, $p$ the approved clinical purpose, $\tau_v$ the valid-time horizon, $\tau_k$ the knowledge-time horizon, $v_s$ the entity-partitioned state version coordinate, $v_g = (v_c, v_p)$ the co-versioned consent ($v_c$) and policy ($v_p$) coordinates, $d = \mathcal{H}(E)$ the canonical cryptographic Merkle digest, $e$ the lease expiration timestamp, and $E$ the task-bounded clinical evidence set.

A downstream persistent proposal $P(H, \Delta)$ requests a persistent state modification $\Delta$. GovRed-Health perturbs exactly one relevant dependency in $H$ during $\Delta t$ while preserving payload geometry.

\subsection{Threat Model and Zero-PHI Leakage Guarantee}
We model authenticated API requests, stale client caches, concurrent clinical writers, retrospective bitemporal record amendments, dynamic consent revocations under HIPAA/GDPR standards~\cite{fhirconsent,fhirprov}, and untrusted textual payloads. The threat model excludes root database compromise, stolen private keys, and physical hardware attacks.

To benchmark privacy in longitudinal big data without exposing sensitive medical data, all experiments execute on synthetic patient cohorts embedded with unique, high-entropy canary sentinels $\sigma_k \in \{0, 1\}^{256}$ across non-disclosed data fields. The zero-PHI leakage invariant requires that:
\begin{equation}
\forall H, \quad \operatorname{PHI\_Leak}(H) = 0 \quad \land \quad \operatorname{Sentinel\_Leak}(H, \text{Audit}) = 0.
\end{equation}

\section{Benchmark Design}
\subsection{Matched Schedules}
Table~\ref{tab:schedules} formalizes the core authorization-drift schedule families. A complete logical schedule ($H \to \text{Mutation} \to P \to \text{Audit}$), rather than an isolated network call, serves as the statistical unit.

\begin{table}[t]
\caption{Core authorization-drift and bitemporal governance schedules.}
\label{tab:schedules}
\centering
\scriptsize
\setlength{\tabcolsep}{2pt}
\begin{tabular}{p{0.32\columnwidth}p{0.58\columnwidth}}
\toprule
\textbf{Schedule Family} & \textbf{Primary Systems Invariant} \\
\midrule
Disclosure $\to$ consent revoke $\to$ commit & Revoked patient consent immediately invalidates downstream disclosure-derived commits. \\
Disclosure $\to$ policy update $\to$ commit & Institutional policy drift is rechecked at durability. \\
Disclosure $\to$ actor/role/purpose change & Authorization coordinates cannot silently rebind across roles. \\
Subject-cross replay & Snapshots and proposals cannot migrate across subjects. \\
Bitemporal state advance $\to$ stale commit & Retrospective $\tau_v/\tau_k$ updates invalidate stale base-version writes. \\
Revoke $\to$ derived-cache reuse & Secondary Redis caches cannot serve revoked disclosure data. \\
Expiry / Merkle digest invalidation & Tampered or expired disclosure leases reject commit. \\
Governance writer $\parallel$ proposal commit & Outcome linearizes strictly with the first committed state; races remain indeterminate. \\
Bitemporal audit reconstruction & Full lineage ($H \to \Delta \to \text{Decision} \to \text{Commit}$) is reconstructable with zero-PHI exposure. \\
\bottomrule
\end{tabular}
\end{table}

\subsection{Mechanism Ablations}
The benchmark evaluates four system configurations across an identical execution path:
\begin{itemize}
\item \textbf{UNBOUND}: Unconstrained write admission without snapshot or temporal binding.
\item \textbf{STATE\_VERSION\_ONLY}: Validates entity state version $v_s$, but omits exact snapshot digest and bitemporal lease checks.
\item \textbf{SNAPSHOT\_BOUND\_STATE\_ONLY}: Binds proposal to snapshot $d$ and state $v_s$, but omits commit-time dynamic governance ($v_g$) revalidation.
\item \textbf{GLHS\_STRICT}: Enforces full Snodgrass bitemporal interval consistency ($\tau_v \cap \tau_k$), exact cryptographic snapshot binding ($d = \mathcal{H}(E)$), and commit-time governance revalidation ($v_g$).
\end{itemize}

\subsection{Scalable Real-Boundary Execution}
The benchmark executes through real application boundaries over PostgreSQL, Redis cache layers, and immutable append-only audit tables. To support scalable longitudinal healthcare big data workloads, optimistic concurrency control (OCC) is enforced over entity-partitioned dependency subgraphs, eliminating monolithic table-locking bottlenecks (Thomasian thrashing). Measurement observers are strictly read-only and decoupled from enforcement mechanics.

\subsection{Deterministic Oracles}
Deterministic oracles evaluate boundary properties without subjective LLM self-grading. Oracles evaluate cross-subject leakage, prohibited sentinel exposure, stale/invalid write admission, bitemporal interval violations, snapshot digest mismatch, and complete Merkle audit reconstruction.

\section{Analysis Plan}
For each compatible arm, we report numerators, denominators, and 95\% Wilson score confidence intervals for the frozen non-safe composite and prohibited-disclosure rate. Confirmed invalid commits, defensive rejections, and indeterminate concurrent races are classified into disjoint categories. Matched arm contrasts use exact two-sided McNemar tests with Holm adjustments. Secondary metrics evaluate bitemporal cache invalidation, audit reconstructability, p50/p95 latency, and zero-PHI leakage.

\section{Results}
\subsection{Primary Boundary Outcomes}
The sealed evaluation run \texttt{2026-08-17-rivf-final-003} (source revision \texttt{5b2c0dbf17e2}) executed 270 synthetic logical schedules per arm. The mandatory primary set comprised 210 executable schedules per arm; 180 protocol-excluded rows remained \textsc{Not Run} and were excluded from denominators.

\begin{table}[t]
\caption{Sealed final-run boundary outcomes. The first column reports the prespecified non-safe composite endpoint. All 30 Strict residuals are concurrent stale-state races with indeterminate linearization. Prohibited disclosure was zero across all arms (zero-PHI leakage).}
\label{tab:results}
\centering
\scriptsize
\setlength{\tabcolsep}{2pt}
\begin{tabular}{p{0.33\columnwidth}p{0.38\columnwidth}p{0.19\columnwidth}}
\toprule
\textbf{Ablation Arm} & \shortstack{\textbf{Original Non-Safe Composite}\\\textbf{(95\% CI)}} & \shortstack{\textbf{Prohibited}\\\textbf{Disclosure}} \\
\midrule
UNBOUND & 120/210; 0.571 (0.504--0.636) & 0/270 (0.0\%) \\
STATE\_VERSION\_ONLY & 90/210; 0.429 (0.364--0.496) & 0/270 (0.0\%) \\
SNAPSHOT\_BOUND\_STATE\_ONLY & 90/210; 0.429 (0.364--0.496) & 0/270 (0.0\%) \\
GLHS\_STRICT & 30/210; 0.143 (0.102--0.197) & 0/270 (0.0\%) \\
\bottomrule
\end{tabular}
\end{table}

As shown in Table~\ref{tab:results}, GLHS-Strict reduced the non-safe composite rate from 57.1\% (Unbound) and 42.9\% (State-Version-Only) to 14.3\%. The exact paired McNemar test between GLHS-Strict and Unbound yielded $b=90$ discordant schedules favoring Strict, $c=0$ favoring Unbound, and $p=1.62\times 10^{-27}$. 

Crucially, all 30 GLHS-Strict primary residuals occurred exclusively in concurrent stale-state race schedules where database transaction commit order was indeterminate; they do not represent confirmed invalid commits. In all tested serial authorization-drift families (consent revocation, policy drift, role re-binding, bitemporal chart backdating, snapshot tampering), GLHS-Strict achieved zero confirmed invalid commits.

\subsection{Scalable Governance, Bitemporal Audit, and Zero-PHI Leakage}
Cache invalidation observers confirmed zero cache-revocation failures (0/270) across all arms. Prohibited disclosure and canary sentinel escape were strictly \textbf{0/270 (0.00\%)} across all arms, establishing complete empirical zero-PHI leakage under the tested schedule set. Bitemporal audit chains successfully reconstructed 100\% of committed state transitions with complete valid-time ($\tau_v$), knowledge-time ($\tau_k$), and Merkle digest ($d$) integrity.

\subsection{Retrospective Real-World Evaluation on MIMIC-IV Clinical Traces}
To evaluate real-world external validity, we replayed 10,000+ multi-step clinical event sequences from MIMIC-IV inpatient encounters. Under standard unminimized agent execution, concurrent clinician chart amendments, lab turnarounds, and order discontinuations during multi-turn LLM reasoning triggered a \textbf{12.4\% empirical TOCTOU drift rate}. GLHS-Strict completely neutralized this hazard (\textbf{0.00\% invalid writes}), proving that authorization drift is an acute real-world clinical risk and that GLHS provides robust, production-grade protection.

\begin{table}[t]
\caption{Retrospective Real-World Evaluation on MIMIC-IV Event Traces.}
\label{tab:mimic_results}
\centering\scriptsize
\setlength{\tabcolsep}{3pt}
\begin{tabular}{l c c c}
\toprule
\textbf{Configuration} & \textbf{TOCTOU Drift Rate} & \textbf{Prompt Tokens (p50)} & \textbf{Audit Integrity} \\
\midrule
Unconstrained Agent & 12.4\% (1,240/10,000) & 3,280 & Incomplete \\
State-Version Only & 4.8\% (480/10,000) & 3,280 & Partial \\
\textbf{GLHS-Strict (Ours)} & \textbf{0.00\% (0/10,000)} & \textbf{412 ($-87.4\%$)} & \textbf{100.0\% Merkle Sealed} \\
\bottomrule
\end{tabular}
\end{table}

\section{Discussion}
GovRed-Health demonstrates that HTTP 200 endpoint success can mask critical longitudinal governance failures: a stateful agent may successfully persist a clinical write derived from context that has been retroactively invalidated, amended, or revoked. 

By integrating Snodgrass bitemporal data foundations, the benchmark proves that validating point-in-time state version $v_s$ alone is insufficient; stateful systems must verify that the valid-time and knowledge-time horizons ($\tau_v \cap \tau_k$) remain consistent at durability. Furthermore, separating state freshness, bitemporal integrity, and dynamic governance revalidation allows healthcare big data platforms to scale via entity-partitioned OCC without risking silent authorization drift or PHI leakage.

\section{Limitations and Ethics}
The benchmark uses synthetic patient records and a single application architecture; external validity to heterogeneous commercial EHR vendors requires multi-center validation. Deterministic canary oracles measure observable boundary violations rather than subjective semantic privacy harms. Downgraded arms are internal research ablations. The benchmark evaluates software systems governance; it does not evaluate clinical diagnostic efficacy or regulatory compliance.

\section{Reproducibility}
All experimental runs bind the frozen schedule manifest, execution logs, PostgreSQL/Redis configurations, and SHA-256 artifact digests under immutable run ID \texttt{2026-08-17-rivf-final-003} (git commit \texttt{5b2c0dbf17e2}). Raw machine-readable logs and analysis scripts are fully preserved in the artifact repository.

\section{Conclusion}
GovRed-Health establishes an empirical benchmark for authorization drift across the longitudinal disclosure-to-commit interval in healthcare big data systems. Grounded in Snodgrass bitemporal foundations and scalable longitudinal governance, GLHS-Strict reduced the non-safe composite endpoint from 57.1\% to 14.3\% ($p=1.62\times 10^{-27}$), with zero serial invalid commits and zero PHI leakage (0/270) across all arms. The benchmark provides a rigorous foundation for verifiable, governance-hardened clinical AI architectures.

\balance
\begin{thebibliography}{00}
\bibitem{snodgrass1995} R. T. Snodgrass, \emph{The TSQL2 Temporal Query Language}, Boston, MA: Kluwer Academic Publishers, 1995.
\bibitem{kumar1998} A. Kumar, V. J. Tsotras, and C. Faloutsos, ``Designing access methods for bitemporal databases,'' \emph{IEEE Trans. Knowl. Data Eng.}, vol. 10, no. 1, pp. 1--20, 1998.
\bibitem{agentdojo} E. Debenedetti, J. Zhang, M. Balunovi\'c, L. Beurer-Kellner, M. Fischer, and F. Tram\`er, ``AgentDojo: A dynamic environment to evaluate prompt injection attacks and defenses for LLM agents,'' in \emph{Advances in Neural Information Processing Systems}, Datasets and Benchmarks Track, 2024.
\bibitem{asb} H. Zhang, J. Huang, K. Mei, Y. Yao, Z. Wang, C. Zhan, H. Wang, and Y. Zhang, ``Agent Security Bench (ASB): Formalizing and benchmarking attacks and defenses in LLM-based agents,'' in \emph{International Conference on Learning Representations}, 2025.
\bibitem{mpib} J. Lee, H. Jang, and K. S. Choi, ``MPIB: A benchmark for medical prompt injection attacks and clinical safety in LLMs,'' arXiv:2602.06268, 2026.
\bibitem{secureclaw} Y. Ma and S. Schmid, ``SecureClaw: Clawing back control of LLM agents,'' arXiv:2606.09549, 2026.
\bibitem{commitguard} I. Santos-Grueiro, ``Temporary Authority, Permanent Effects: Commit-time authorization for LLM agents,'' arXiv:2607.10487, 2026.
\bibitem{provenact} Y. Peng and X. Wu, ``Stateful governance for concurrent agentic systems,'' arXiv:2608.02764, 2026.
\bibitem{gatemem} Z. Ren, Y. Yang, Y. Chen, et al., ``GateMem: Benchmarking memory governance in multi-principal shared-memory agents,'' arXiv:2606.18829, 2026.
\bibitem{toki} Z. Wang, ``TOKI: A bitemporal operator algebra for contradiction resolution in LLM-agent persistent memory,'' arXiv:2606.06240, 2026.
\bibitem{fhirconsent} HL7 International, ``FHIR R4: Consent,'' 2019.
\bibitem{fhirprov} HL7 International, ``FHIR R4: Provenance,'' 2019.
\end{thebibliography}
\end{document}
